\section{Related Work}\label{sec:related}

\subsection{Short-Term Traffic Speed and Congestion Forecasting}
Short-term traffic forecasting methods span three major paradigms: classical statistical models, tabular machine learning, and deep spatiotemporal architectures. Classical approaches such as Historical Average (HA), ARIMA, and Support Vector Regression (SVR) offer transparent behavior and low computational overhead, making them essential baselines, but they struggle to capture non-linear congestion transitions under abrupt demand surges.

Deep spatiotemporal networks have achieved strong benchmark performance by modeling road networks as graphs. Diffusion Convolutional Recurrent Neural Networks (DCRNN)~\cite{li2018dcrnn} capture bidirectional spatial diffusion alongside recurrent temporal dynamics; Spatio-Temporal Graph Convolutional Networks (STGCN)~\cite{yu2018stgcn} replace recurrence with gated temporal convolutions for faster training; and Graph WaveNet~\cite{wu2019graphwavenet} introduces adaptive adjacency matrices and dilated causal convolutions. While effective on dense sensor benchmarks, these models act as opaque black boxes, require GPU infrastructure, and rarely report event-level early-warning accuracy under controlled alert budgets.

For structured tabular telemetry with moderate spatial dimensionality (such as a 10-segment corridor), histogram-based gradient boosted decision trees---popularized by LightGBM~\cite{ke2017lightgbm}---offer competitive accuracy, rapid CPU training, native handling of missing values and categorical segment identifiers, and straightforward reproducibility. Furthermore, raw classifier outputs in rare-event regimes are frequently miscalibrated~\cite{guo2017calibration}; non-parametric isotonic regression~\cite{zadrozny2002isotonic} provides a principled post-hoc calibration mechanism when fitted on held-out chronological splits.

\subsection{LLMs and Tool-Augmented Agents in Transportation}
Recent research has explored integrating LLMs into intelligent transportation systems (ITS) along two directions. First, foundation-model adaptations such as ST-LLM~\cite{liu2024stllm} and xTP-LLM~\cite{guo2024xtpllm} align spatiotemporal traffic embeddings with partially frozen LLMs to produce forecasts and textual explanations. Second, agentic frameworks inspired by ReAct~\cite{yao2023react}---notably TrafficGPT~\cite{zhang2024trafficgpt}---connect LLMs to external traffic simulation and analysis tools for interactive querying and decision support.

Despite their conversational flexibility, existing LLM traffic studies exhibit three recurring limitations for operational traveler advisories: (1) reliance on offline benchmark datasets or microscopic simulators rather than live open roadside feeds with real-world missingness and latency; (2) insufficient separation between numerical computation and text generation, leaving responses vulnerable to subtle numerical hallucination; and (3) omission of verified domain regulations (e.g., holiday HOV controls and ramp closures) and explicit scope-refusal boundaries.

\subsection{Research Gap and Positioning}
Table~\ref{tab:mae} and Section~\ref{sec:method} reflect how our system bridges these gaps: instead of asking an LLM to predict traffic directly or relying on an opaque black-box model alone, we pair lightweight tabular forecasting and isotonic calibration with a verified regulatory knowledge base, live incident feeds, and deterministic verification guardrails. Crucially, we evaluate both predictive utility (against persistence and historical average under strict chronological splits) and agent reliability (refusal accuracy, numerical fidelity, and fallback behavior).
